\section{Threats to validity\label{sec:threats}}

\paragraph{Construct validity.}
Junior, Mid, and Senior are unsupervised clusters based on observable GitHub activity, not verified job titles or measured skills. The label reflects relative activity within the analyzed population, not career stage. Three of the twelve clustering indicators - Agentic-PR activity span, agent diversity, and Agentic-PR repository range - directly describe agent-related activity, rather than independent, pre-experiment experience without the help of AI tools. The profiles are therefore partly built from the same kind of behavior that RQ1 and RQ2 later measure. The remaining indicators describe general GitHub activity unrelated to agent use. Bot and agent detection relies on AIDev's own labels and simple login-pattern filtering, which can misclassify unusually named bot accounts or human accounts. In addition, we initially used a bot/human flag collected from GitHub's GraphQL API - it turned out to be quite unreliable, and in most cases reported a user as human even when the profile description suggested otherwise. Furthermore, the profile-enrichment data was collected about 9 months after the AIDev cutoff date - indicators such as the number of followers or yearly contribution counts therefore reflect the account's state at the later data-collection date, not at the time of its Agentic-PR, so part of the activity profile may come from a period later than the contribution it is meant to characterize.

\paragraph{Internal validity.}
The study design is observational: none of the three research questions supports a causal conclusion about the effect of experience or of agent use. Pull-request complexity, task type, and repository governance norms are not controlled and likely confound the acceptance rate. The RQ2 finding that acceptance decreases with activity level, for example, is equally consistent with the hypothesis that more active developers take on harder tasks, and with the hypothesis that agents serve them less well. The parameter $k=3$ was kept for interpretability, even though the statistics do not point to a strict numeric optimum (Section 3). A different choice of $k$, or a different clustering algorithm, could give a different group composition and, as a result, a different comparison of outcomes. The 1755 developers for whom GraphQL indicators were not collected were estimated from features with weak explanatory power (observed $R^2$ between $-0.27$ and $0.10$, Section 3) - so their assignment to profiles is relatively uncertain.

A complete-case sensitivity check - clustering only the 70,434 developers with real, non-estimated data, instead of the full 72,189 - initially showed that the version of the pipeline from Section 3 based on the quantile transformation was unstable under this change in population: the cluster assignment agreed with clustering on the full population, for the same overlapping developers, at only an Adjusted Rand Index of 0.46, and the Junior vs Mid ordering in RQ1 and RQ2 reversed. We traced this to the quantile-transformation step: because several indicators are heavily zero-inflated, the quantile transformation maps values based on their rank within the current sample, so removing even a small part of the population shifts these ranks for everyone else. Repeating the same complete-case comparison with the log(1+x) transformation plus standardization gave an Adjusted Rand Index of 0.98, a practically stable result, which is why this method was used in preprocessing.

We also re-ran RQ1 and RQ2, restricted to the 70,434 developers, under the now-adopted log(1+x) pipeline (RQ3 does not depend on clustering and is not affected by this). Both results were essentially unchanged (e.g., RQ1: Junior/Mid/Senior means of 3.59/42.27/7.18 vs.\ 3.57/42.12/7.94 in the full-population analysis; RQ2: 90.3\%/91.8\%/83.4\% vs.\ 90.2\%/91.8\%/83.2\%), with the same pattern of significance in the pairwise comparisons in both cases. For transparency: a third candidate outcome, measuring the external comment count per Agentic-PR across the same three profiles, was explored during work on the paper, but excluded from this study after it failed the same check - its direction reversed, and statistical significance disappeared entirely under the complete-case restriction, and it was also based on a small and unevenly distributed sample.

In addition, we ran four further sensitivity checks against the log(1+x)-based pipeline.

1) Alternative k: for k = 4 and 5, the three main clusters stay close to their sizes at k = 3, with the extra clusters splitting off from Senior as small, higher-activity fragments, rather than through a thorough reorganization.

2) Repeated initializations: re-running the whole pipeline (Isolation Forest and k-means) with five different random seeds gave an Adjusted Rand Index between 0.98 and 0.99 relative to the reported clustering, showing very high stability.

3) Kept outliers: skipping Isolation Forest entirely gives an Adjusted Rand Index of 0.64 relative to the reported clustering, driven almost entirely by the Senior profile, which shrinks from 13,908 to 1586 developers when the most extreme profiles are isolated into their own cluster instead of being removed beforehand. Junior and Mid stay relatively stable.

4) Excluding the AIDev-derived indicators (Agentic-PR activity span, agent diversity, and repository range), leaving only the remaining nine indicators: this is the least stable of the four checks, with an Adjusted Rand Index of only 0.32, and cluster sizes changing substantially (Mid grows from 12,705 to 28,329). Re-running RQ1 and RQ2 with this 9-indicator clustering: the qualitative RQ2 result holds - Senior still has the lowest acceptance rate (80.7\% vs.\ 90.4\% for Junior and 90.2\% for Mid) - but the Mid vs Senior contrast in RQ1 is much weaker (mean Agentic-PRs 17.6 vs.\ 20.8, close to reversing) than under the full 12-indicator clustering, which is consistent with the construct-validity concern raised above - part of the large RQ1 effect size comes from the fact that the clustering itself is partly built from agent-use indicators.

\paragraph{External validity.}
The dataset covers Agentic-PRs from five specific coding agents in public GitHub repositories up to 1 August 2025, in an early period of agent adoption. Usage patterns, acceptance norms, and maintainer behavior toward agent-generated code may change as the tools and community norms mature, so the reported relationships should not be carried over to later periods or to private/company repositories without re-checking them. RQ3 is further limited to a baseline of repositories with more than 500 stars, because the human-authored comparison sample was collected by the AIDev authors only for repositories above this popularity threshold; this result specifically describes high-visibility, high-star repositories and may not generalize to the long tail of smaller projects that make up most of the full population.

\paragraph{Conclusion validity.}
The distributions of the outcome variables are strongly non-normal and zero-inflated, which justified the use of non-parametric tests and the Bonferroni correction throughout this work (Section 3). With sample sizes in the tens of thousands, even very small differences can reach statistical significance regardless of practical importance; this is most visible for RQ2 ($\varepsilon^2=0.043$), whose effect size is small even though its $p$-value is not. RQ1 is an exception: its effect size ($\varepsilon^2=0.32$) is large by conventional standards, so its significance reflects a meaningful relationship, not just a large sample. RQ3 further rests on clearly unequal group sizes (56,387 Agentic-PRs vs.\ 2,286 human-authored pull requests), which further limits how much weight can be placed on its point estimate, especially for the smaller group.
